\section{Related Work}

\subsection{Video Diffusion Models}
The diffusion and flow-matching frameworks~\cite{lipman2022flow,esser2024scaling} have rapidly advanced from high-fidelity image synthesis to temporally coherent video generation. Early video diffusion methods extended image architectures by interleaving temporal attention or 3D convolution layers to capture inter-frame dependencies~\cite{blattmann2023stable,guo2023animatediff,chen2024videocrafter2}, while later systems scaled this paradigm to photorealistic, multi-second outputs through large-scale pretraining on diverse video corpora~\cite{sora,kong2024hunyuanvideo,yang2024cogvideox,xu2024easyanimate,allegro2024,fan2025vchitect}. A common strategy is to operate within a compressed VAE latent space, which significantly reduces the computational cost of the denoising process and enables generation at higher resolutions and longer durations~\cite{blattmann2023stable,hacohen2024ltx,wan2025wan}. More recently, Diffusion Transformers~\cite{peebles2023scalable} have become the dominant backbone, and various autoregressive and streaming extensions have been proposed to generate longer sequences~\cite{chen2024diffusion,henschel2024streamingt2v,gu2025long,chen2025skyreels,xie2024progressive}. Despite these advances, most video diffusion models treat the generation process as fundamentally two-dimensional, in that frames are synthesized sequentially or in parallel without an explicit model of the underlying 3D scene geometry. As a consequence, generated videos frequently exhibit geometric drift, parallax violations, and inconsistent scene structure when subjected to large camera motions or extended temporal horizons.

\subsection{Camera-controllable Video Generation}
A parallel line of work seeks to condition video diffusion models on explicit camera trajectories to enable controllable viewpoint changes. Representative methods inject camera pose information through additional control modules~\cite{he2024cameractrl,he2025cameractrl,feng2024i2vcontrol,zhang2026panflow}, epipolar attention mechanisms~\cite{yu2024viewcrafter,zhou2025stable}, or 3D-aware rendering signals fed as conditioning inputs~\cite{gu2025das,ren2025gen3c,yang2025omnicam}. While these approaches offer fine-grained camera control within a single generation pass, they do not maintain a persistent scene representation across generation steps. Each clip is generated independently of previously synthesized content, so revisiting a region or extending a trajectory beyond the context window of the model may introduce inconsistencies. The absence of a shared spatial memory limits their applicability to long-horizon, exploratory world generation where 3D coherence must be preserved across many sequential generation rounds.

\subsection{Spatial Memory for Video Generation}
To maintain spatial consistency, a growing body of work augments video diffusion pipelines with explicit spatial memory structures that maintain a 3D scene representation across generation steps. A representative strategy is to lift observed or generated RGB frames into point clouds using estimated depth, accumulate them into a persistent 3D cache, and render target-view conditioning images at each step~\cite{yu2024wonderjourney,engstler2024invisible,yu2024wonderworld,huang2025voyager,zhao2026spatia,chen2025flexworld}. More recent systems further incorporate long-term context retrieval~\cite{yu2025context,xiao2025worldmem} or surfel-indexed view memory~\cite{li2025vmem} to improve temporal consistency over extended sequences. Such memory-based approaches have demonstrated clear improvements in multi-view coherence, since the 3D cache anchors each new frame to a shared geometric scaffold. Nevertheless, existing spatial memory designs operate entirely in RGB pixel space, which is both computation-intensive and is vulnerable to accumulated errors in repeated encoding-decoding operations. These limitations motivate our approach, which constructs and queries spatial memory entirely within the latent space of the diffusion model, preserving representational fidelity while eliminating the rendering and re-encoding overhead that constrains prior methods.
